% ============================================================
%  TEMA 3 — APARTADO 3: EQUILIBRIO ENTRE OFERTA Y DEMANDA DE RRHH
%
%  RECORDATORIO AL EQUIPO:
%  Tras comparar la oferta y la demanda de RRHH se pueden dar tres situaciones:
%
%  1. OFERTA = DEMANDA → Equilibrio. Situación ideal pero poco frecuente.
%     Acciones: mantenimiento, formación para adaptarse a cambios menores.
%
%  2. DEMANDA > OFERTA → Déficit de personal. Necesidad de incorporar.
%     Acciones: reclutamiento externo, promoción interna, horas extra,
%     trabajo temporal, externalización (outsourcing), reincorporación de
%     exempleados (boomerang employees), automatización de tareas para
%     reducir necesidades.
%
%  3. OFERTA > DEMANDA → Superávit de personal. Necesidad de reducir.
%     Acciones: ERE, ERTE, jubilaciones anticipadas, congelación de
%     contrataciones, reducción de jornada, recolocación (outplacement),
%     formación para reasignar a otras áreas. (Se desarrolla en el
%     apartado de procesos sustractivos, sección 10)
%
%  PARA ESTA SECCIÓN INDICAD:
%  - ¿Cuál es la situación de CaixaBank actualmente (déficit/equilibrio/superávit)?
%  - ¿Ha variado en los últimos años (especialmente tras la fusión)?
%  - ¿Qué acciones concretas ha tomado para lograr el equilibrio?
%  - ¿Está satisfecha con los resultados de esas acciones?
% ============================================================

\section{Equilibrio entre oferta y demanda de recursos humanos}

\subsection{Situación actual de CaixaBank}

% TODO: Tras el ERE de 2021 (fusión con Bankia, con salida de ~6.000 empleados),
% la plantilla se ha ido ajustando. Pero al mismo tiempo la digitalización
% crea demanda de nuevos perfiles tecnológicos que no siempre se pueden cubrir
% internamente. Puede haber simultáneamente superávit en ciertos perfiles
% tradicionales (cajeros, operadores de red) y déficit en perfiles digitales.

[Descripción a completar]

\begin{table}[H]
  \centering
  \small
  \caption{Situación de equilibrio oferta-demanda en CaixaBank por tipo de perfil}
  \begin{tabularx}{\textwidth}{p{0.30\textwidth} p{0.20\textwidth} X}
    \toprule
    \textbf{Familia de puestos}    & \textbf{Situación} & \textbf{Acciones adoptadas} \\
    \midrule
    Perfiles bancarios tradicionales (oficina) & [Superávit/Déficit/Equilibrio] & \\
    Perfiles tecnológicos (datos, ciberseg., digital) & [Superávit/Déficit/Equilibrio] & \\
    Perfiles de gestión y dirección  & [Superávit/Déficit/Equilibrio] & \\
    Perfiles de cumplimiento y riesgo & [Superávit/Déficit/Equilibrio] & \\
    \bottomrule
  \end{tabularx}
\end{table}

\subsection{Acciones de ajuste ante déficit de personal}

% TODO: En qué áreas hay escasez y cómo la cubre CaixaBank.
% ¿Recurre a reclutamiento externo? ¿A la promoción interna? ¿A la formación
% de recualificación (reskilling)?

\begin{itemize}
  \item \textbf{Reclutamiento externo:} [descripción]
  \item \textbf{Promoción interna:} [descripción]
  \item \textbf{Reskilling / upskilling:} [¿tienen programas de recualificación?]
  \item \textbf{Trabajo temporal o externalización:} [descripción]
  \item \textbf{Acuerdos con universidades / campus recruitment:} [descripción]
\end{itemize}

\subsection{Acciones de ajuste ante superávit de personal}

% Se desarrollará con mayor detalle en la sección de Procesos Sustractivos (sección 10).
% Aquí haced un resumen de los mecanismos de reducción utilizados.

\begin{itemize}
  \item \textbf{ERE 2021 (fusión Bankia):} [número de afectados, condiciones]
  \item \textbf{Jubilaciones anticipadas:} [descripción]
  \item \textbf{Recolocación interna (movilidad geográfica / funcional):} [descripción]
  \item \textbf{Planes de formación para reconversión profesional:} [descripción]
\end{itemize}

\begin{notaequipo}
  Preguntad en la entrevista:
  \begin{enumerate}
    \item ¿En qué perfiles sienten más escasez de candidatos en el mercado?
    \item ¿Cómo afrontaron el exceso de plantilla tras la fusión con Bankia?
    \item ¿Han implementado programas de reskilling o upskilling para
          reconvertir perfiles tradicionales en perfiles digitales?
    \item ¿Cómo gestionan las situaciones de superávit temporal
          sin recurrir a despidos (ERTE, reducción de jornada...)?
    \item ¿Está la plantilla actual equilibrada o hay áreas con tensión?
  \end{enumerate}
\end{notaequipo}
